\subsection{Institutions in 1924 and the age target of the 1931 plan}
Table XVII of the Central Statistical Administration's survey of educational
institutions on 1 January 1924 reports literacy schools and points, teachers
and pupils by sex, separately for rural areas. I transcribe 30 rows for
15 territories from printed pages 462--467 of \emph{Narodnoe obrazovanie v
SSSR}, vol. XXVIII, issue 1 (1926). General adult schools in adjacent
columns are excluded. Fourteen territories can be linked to validated male
and female cells aged 18--22 in 1926. The Bashkir republic lacks this
validated cohort in the estimation sample. No duplicated oblast exposure
is assigned to the individual Ural districts.

The counts describe institutions and pupils at a date, rather than unique
graduates or exposure of a particular age cohort. Coverage is heterogeneous:
only five of 23 rural Mari institutions were surveyed, compared with all
81 in the Tatar republic. The total Bashkir teacher count exceeds the sum
of the two reported sexes by three; the transcription preserves this
unresolved discrepancy. Territorial transfers between 1924, 1926 and 1939
have not been fully harmonised.

Exploratory models relate the subsequent female-minus-male literacy gain
to female shares of surveyed pupils or teachers, controlling for the
initial literacy gap. Other models use log counts of institutions or
teachers, with cohort population as a territory-size control. This
population is not treated as the full eligible population or used as an
enrolment-rate denominator. Across 24 combinations of exposures, weights
and coverage restrictions, every CR2/Satterthwaite 95 per cent interval
includes zero. Results do not establish a robust institutional mechanism;
small samples, incomplete observation and endogenous allocation also
prevent interpreting them as evidence of no effect.

The 1931 national economic plan specifies ages 16--50 in towns and 16--40
in rural areas, with some instruction above 40 in collectivised areas
(printed pages 67--68). Planned contingents are not observed graduates.
As a sensitivity check, dropping the cohort aged 33--37 in 1926, whose
ages straddle 40 in 1931, changes the main female contrast from \SN{policy.baseline.female.b} to
\SN{policy.omit_33_37.female.b} percentage points. The younger group remains aged 18--32 in both models; ages 33--37 belong to the intermediate group. After their omission, the female contrast has a CR2 interval of
[\SN{policy.omit_33_37.female.ci_low_cr2},\SN{policy.omit_33_37.female.ci_high_cr2}]. The triple difference changes from \SN{policy.baseline.DDD.b} to \SN{policy.omit_33_37.DDD.b}.
This checks the inclusion of one intermediate cohort; it does not establish an untreated older
comparison group or remove migration and reporting bias.

Machine-readable transcriptions, source page images, all specifications
and fixed seeds are supplied in \texttt{supplementary/likbez\_1924\_analysis/}.
The full source is available at
\url{https://istmat.org/files/uploads/22102/obrazovanie_1921-25_vyp.1_1924_god.pdf};
the 1931 plan at
\url{https://istmat.org/files/uploads/52124/plan_sssr_na_1931_god.pdf}.
